% Capitulo: Construir a própria aplicação de master data e governança
\chapter{Construir a própria aplicação de master data e governança}
\label{cap:aplicacao-propria}

Há muito material sobre ferramentas de mercado para master data e governança, e muito
pouco sobre o que considerar quando uma empresa decide construir a sua. Até pouco tempo,
fazia sentido: construir era caro demais para a maioria. Com a IA ajudando a escrever
código, isso mudou, e mais times vão se perguntar se vale a pena. Este capítulo é um guia
para essa decisão e para a construção. Não é um tutorial de uma tecnologia específica; é a
lista de decisões que vale tomar antes de escrever a primeira linha de código.

\section{Construir ou comprar}
\label{sec:construir-comprar}

A IA baixou muito o custo de construir. Não baixou o custo de manter. Uma aplicação interna
precisa de alguém que corrija erros, atualize dependências, atenda dúvidas e adapte o
sistema quando o negócio mudar. A pergunta que decide entre construir e comprar, portanto,
não é \enquote{conseguimos fazer?}, e sim \enquote{quem vai cuidar disso daqui a dois
anos?}. A Tabela~\ref{tab:construir-comprar} reúne os outros critérios.

\begin{table}[htbp]
  \centering
  \caption{Quando construir e quando comprar uma solução de master data e governança.}
  \label{tab:construir-comprar}
  \small
  \begin{tabularx}{\textwidth}{@{}>{\raggedright\arraybackslash}p{0.17\textwidth}
      >{\raggedright\arraybackslash}X>{\raggedright\arraybackslash}X@{}}
    \toprule
    Critério & Construir faz sentido quando & Comprar faz sentido quando \\
    \midrule
    Escopo & Um ou dois domínios, com regras muito próprias do negócio & Muitos domínios e fontes, com regras comuns de mercado \\
    \addlinespace
    Equipe & Há quem programe e quem mantenha depois & Não há quem mantenha código interno \\
    \addlinespace
    Regulação & Um histórico de mudanças simples basta & Auditoria formal, certificações, setor regulado \\
    \addlinespace
    Prazo & É preciso algo em semanas, um domínio por vez & É um programa corporativo de longo prazo \\
    \addlinespace
    Custo & A licença é desproporcional ao problema & Manter código próprio custaria mais que a licença \\
    \addlinespace
    Integração & É preciso encaixar em fluxos internos específicos & A ferramenta já conecta com toda a pilha de dados \\
    \bottomrule
  \end{tabularx}
  \fonte{elaborado pelo autor.}
\end{table}

Construir também não precisa ser para sempre. Começar com uma aplicação própria num
domínio, aprender com ela o que a empresa realmente precisa e migrar para uma ferramenta de
mercado quando a escala pedir é um caminho legítimo. A condição é desenhar a aplicação de
forma que os dados saiam fácil: registro mestre, tabela de correspondência e histórico de
decisões em tabelas simples, num banco de dados comum, e não presos à lógica da aplicação.

\section{Começar pelo problema, não pela plataforma}
\label{sec:comecar-pequeno}

O erro mais comum de quem decide construir é começar pela plataforma: um sistema genérico de
master data, que serviria para qualquer domínio, com telas de configuração e regras
parametrizáveis. Esse caminho reproduz o programa eterno que o capítulo~\ref{cap:master-data}
criticou, agora com código próprio.

O ponto de partida melhor é um problema concreto, com um domínio, uma regra e um dono.
Algo como: \enquote{clientes com nomes diferentes no ERP e no CRM, revisados pelo time
comercial, para o projeto de acompanhamento de carteira}. A plataforma, se vier, nasce do
segundo e do terceiro problema, quando os padrões começam a se repetir.

O primeiro usuário da aplicação é o projeto de negócio que puxou o domínio
(\veja{sec:master-data-puxado}). É ele quem diz se a aplicação está resolvendo, e é a
opinião dele que deve orientar as primeiras melhorias, não uma lista de recursos imaginada
de antemão.

\section{Uma arquitetura mínima}
\label{sec:arquitetura-minima}

Uma aplicação de master data e governança, mesmo pequena, tem algumas peças que não podem
faltar (Figura~\ref{fig:aplicacao-propria}).

\begin{figure}[htbp]
  \centering
  % Arquitetura minima de uma aplicacao propria de master data e governanca (cap. 15).
\begin{tikzpicture}[diagrama, x=1mm, y=1mm]
  % fontes e ingestao
  \foreach \y/\nome/\id in {8/ERP/f1, 0/CRM/f2, -8/Planilhas/f3}
    \node[caixa, minimum width=14mm, minimum height=6mm] (\id) at (7,\y) {\nome};
  \node[moldura, fit=(f1)(f3), label={[rotulo]above:Fontes}] {};
  \node[caixa, minimum width=17mm, minimum height=10mm] (ing) at (29,0) {Ingestão\\{\scriptsize (prata)}};
  \foreach \id in {f1,f2,f3} \draw[seta] (\id.east) -- (ing.west);

  % motor de conciliacao
  \node[destaquecaixa, minimum width=24mm, minimum height=18mm] (mot) at (54,0)
    {Motor de\\conciliação\\[0.5mm]{\scriptsize regras,}\\{\scriptsize similaridade, IA}};
  \draw[seta] (ing) -- (mot);

  % saidas
  \node[bancodestaque, minimum width=24mm, minimum height=13mm] (reg) at (102,9)
    {\scriptsize Registro mestre e\\\scriptsize correspondência};
  \node[caixa, minimum width=20mm, minimum height=9mm] (fila) at (81,-10) {Fila de\\revisão};
  \node[caixa, minimum width=18mm, minimum height=9mm] (stew) at (104,-10) {Steward\\revisa};
  \node[caixa, draw=cinza, dashed, text=cinza, minimum width=20mm, minimum height=8mm] (notif) at (81,-24)
    {\scriptsize Mensagem\\\scriptsize aos donos};
  \draw[seta] (mot.east) ++(0,3) -- node[rotulo, above, sloped] {alta confiança} (reg.west);
  \draw[seta] (mot.east) ++(0,-3) -- node[rotulo, pos=0.55, below left=0.5mm] {dúvida} (fila.west);
  \draw[seta] (fila) -- (stew);
  \draw[seta] (stew) -- (reg);
  \draw[seta, draw=cinza] (fila) -- (notif);
  \node[rotulo, anchor=south] at (reg.north) {lido por BI, sistemas e IA};

  % aprendizado
  \draw[seta, dashed, draw=cinza] (stew.south) -- ++(0,-19) -| (mot.south);
  \node[rotulo, anchor=north] at (68,-33) {as decisões do steward viram regras e exemplos para o motor};
\end{tikzpicture}

  \caption{Uma arquitetura mínima para uma aplicação própria de master data e
    governança.}
  \label{fig:aplicacao-propria}
  \fonte{elaborado pelo autor.}
\end{figure}

A \textbf{ingestão} traz os registros das fontes. Sempre que possível, ela deve ler da
camada prata (capítulo~\ref{cap:medalhao}), onde os dados já estão limpos e tipados, em vez
de ir direto aos sistemas de origem. O \textbf{motor de conciliação} compara os registros e
atribui a cada par uma nota de confiança. A partir da nota, cada par segue um de três
caminhos: com confiança alta, vai direto para o registro mestre; com confiança baixa, é
descartado; no meio, vai para a \textbf{fila de revisão}, onde uma pessoa decide. O
\textbf{registro mestre}, com a tabela de correspondência, é o produto final, lido pelo BI,
pelos sistemas e pela IA. As \textbf{notificações} avisam quem precisa decidir ou saber.

A seta tracejada na parte de baixo da figura é a que transforma uma ferramenta em governança:
as decisões do steward voltam para o motor, como regras novas ou exemplos, e a próxima
rodada sugere melhor. É o ciclo da Figura~\ref{fig:governanca-ciclo} implementado em
código.

Para a tecnologia, a recomendação é usar o mínimo que o time já domina. Um banco de dados,
Python, um agendador que rode o processo depois de cada carga e um canal de mensagens que as
pessoas já usem resolvem o começo. Complexidade adicional só quando um problema real pedir.

\section{O motor de conciliação}
\label{sec:motor-conciliacao}

O motor de conciliação funciona melhor em camadas, da comparação mais barata para a mais
cara. Primeiro, identificadores fortes: se dois cadastros têm o mesmo documento de
registro da empresa, ou o mesmo código de produto, são a mesma entidade, sem discussão.
Depois, regras de normalização e similaridade de texto para o que sobrou. Por fim, só para os
casos que as regras não resolvem, uma sugestão da IA, sempre acompanhada da explicação do
porquê.

A base teórica dessa comparação é a conciliação probabilística de registros
\parencite{fellegi1969,christen2012}, e há bibliotecas de código aberto maduras que a
implementam em escala, como o Splink \parencite{linacre2022}, além de bibliotecas rápidas de
similaridade de texto. Para começar, porém, a biblioteca padrão do Python é suficiente. O
trecho abaixo normaliza os nomes (caixa baixa, sem acento, sem pontuação, sem sufixos
societários) e compara cada par pela similaridade dos textos.

\begin{python}
import unicodedata
from difflib import SequenceMatcher
from itertools import combinations

SUFIXOS = {"ltda", "sa", "me", "eireli"}

def normalizar(nome: str) -> str:
    """Minúsculas, sem acento, pontuação ou sufixo."""
    nome = unicodedata.normalize("NFKD", nome)
    nome = nome.encode("ascii", "ignore").decode().lower()
    nome = "".join(c for c in nome if c.isalnum() or c == " ")
    partes = nome.split()
    return " ".join(p for p in partes if p not in SUFIXOS)

def pares_suspeitos(nomes, limiar=0.85):
    """Pares de nomes que parecem a mesma entidade."""
    for a, b in combinations(sorted(set(nomes)), 2):
        na, nb = normalizar(a), normalizar(b)
        nota = SequenceMatcher(None, na, nb).ratio()
        if nota >= limiar:
            yield a, b, round(nota, 2)

clientes = ["AGRO SILVA LTDA", "Agro Silva", "Agrosilva",
            "Cerealista Norte S.A."]
for a, b, nota in pares_suspeitos(clientes):
    print(f"Revisar: {a!r} x {b!r} ({nota})")
\end{python}

\begin{terminal}
Revisar: 'AGRO SILVA LTDA' x 'Agro Silva' (1.0)
Revisar: 'AGRO SILVA LTDA' x 'Agrosilva' (0.95)
Revisar: 'Agro Silva' x 'Agrosilva' (0.95)
\end{terminal}

O resultado encontra as três grafias do mesmo cliente e ignora a empresa diferente. Num
sistema real, cada par deixaria de ser impresso na tela e viraria um item na fila de revisão,
com uma mensagem para o dono do domínio. Duas observações importam mais do que o código. A
primeira: comparar todos os pares entre si cresce rápido demais quando o cadastro tem
milhares de registros, e é por isso que as técnicas de bloqueio, que só comparam registros
com alguma característica em comum, existem \parencite{christen2012}. A segunda: o limiar
de 0,85 é arbitrário. Ele deve ser calibrado com uma amostra revisada pelo negócio e
ajustado ao longo do tempo com as decisões do steward.

\section{Revisão humana e notificações}
\label{sec:revisao-notificacao}

A fila de revisão é o coração da aplicação. É ali que a governança acontece, e por isso ela
precisa ser fácil de usar. Cada item deve mostrar os dois registros lado a lado, a nota de
confiança, o motivo da sugestão e três ações: aprovar, corrigir ou rejeitar. Uma interface
simples basta, e existem bibliotecas em Python que permitem montar telas internas desse tipo
em poucas horas.

As notificações devem ir para quem precisa decidir, pelo canal que a pessoa já usa (e-mail
ou o chat da empresa), com o link direto para a revisão. Um cuidado: mensagem demais vira
ruído e passa a ser ignorada. Agrupar as pendências por domínio, num resumo diário ou
semanal, costuma funcionar melhor do que avisar cada caso.

Por fim, a revisão precisa de um prazo combinado e de um indicador de itens parados. Uma
fila que ninguém revisa é pior do que nenhuma fila, porque dá a impressão de que o problema
está sendo cuidado.

\section{Histórico, permissões e privacidade}
\label{sec:auditoria-privacidade}

Toda mudança no registro mestre precisa ficar registrada: o que mudou, quem decidiu,
quando e a partir de qual sugestão. Sem esse histórico, a aplicação é só um banco de dados
editável, e não há governança. Com ele, é possível responder a perguntas que sempre
aparecem, como \enquote{por que este cliente foi unido àquele?}, e desfazer uma decisão
errada sem perder o resto.

As permissões devem ser definidas por domínio: quem vê, quem pode sugerir, quem pode
aprovar. O time comercial aprova clientes; o time de compras aprova fornecedores. E, como
cadastros costumam ter dados pessoais, o que for enviado a um modelo de IA externo precisa
ser limpo antes: nomes de pessoas, documentos e contatos não têm por que sair do ambiente da
empresa para que a IA sugira se duas razões sociais são a mesma (\veja{sec:seguranca}).

\section{Operar como produto}
\label{sec:operar-produto}

A aplicação de master data é um produto de dados como qualquer outro, e as práticas da
Parte~\ref{parte:devops} valem para ela: código versionado, testes, entrega contínua e
monitoramento. Ela também precisa de métricas próprias, que mostram se está funcionando:

\begin{itemize}
  \item porcentagem de pares conciliados automaticamente;
  \item tempo médio entre a sugestão e a decisão;
  \item número de itens parados na fila e há quanto tempo;
  \item duplicidades encontradas por carga, que devem cair com o tempo;
  \item custo por execução, inclusive em tokens, se houver IA no motor
    (capítulo~\ref{cap:custo-consumo}).
\end{itemize}

E precisa de um dono, separado dos donos dos domínios. O dono do domínio de clientes decide
quem é cliente; o dono da aplicação garante que ela funciona. Misturar os dois papéis é um
jeito rápido de ter uma aplicação que ninguém mantém.

\rascunho{
  Relato opcional: uma aplicação interna desse tipo que eu construí ou acompanhei, o que
  funcionou e o que eu faria diferente.
}

\section{Lista de verificação}
\label{sec:checklist-aplicacao}

A Tabela~\ref{tab:checklist-aplicacao} reúne, como referência rápida, as perguntas que vale
responder antes de começar. Nenhuma delas é técnica no sentido estrito, e isso diz muito
sobre onde estão os riscos de um projeto desses.

\begin{table}[htbp]
  \centering
  \caption{Perguntas para responder antes de construir a própria aplicação.}
  \label{tab:checklist-aplicacao}
  \small
  \begin{tabularx}{\textwidth}{@{}>{\raggedright\arraybackslash}p{0.22\textwidth}
      >{\raggedright\arraybackslash}X@{}}
    \toprule
    Tema & Pergunta \\
    \midrule
    Domínio e dono & Qual domínio vem primeiro, e quem decide as regras dele? \\
    Fontes & De onde vêm os registros, e com que frequência mudam? \\
    Identidade & O que faz dois registros serem o mesmo? Existe um identificador forte? \\
    Limiares & Acima de que nota aprova sozinho? Abaixo de qual descarta? \\
    Revisão & Quem revisa, em quanto tempo e por qual canal? \\
    Notificação & Quem precisa saber de cada mudança, e com que frequência? \\
    Histórico & Toda mudança fica registrada, com quem decidiu e quando? \\
    Privacidade & O que pode ir para um modelo de IA, e o que precisa ser mascarado? \\
    Testes & Unicidade, correspondência e regras de negócio rodam a cada carga? \\
    Custo & Quanto custa cada execução, inclusive em tokens? \\
    Manutenção & Quem cuida da aplicação daqui a dois anos? \\
    Saída & Se um dia migrar para uma ferramenta de mercado, os dados saem fácil? \\
    \bottomrule
  \end{tabularx}
  \fonte{elaborado pelo autor.}
\end{table}

\begin{leituras}
  \leitura{christen2012}{Para aprofundar o motor de conciliação: normalização,
    bloqueio, comparação e avaliação dos pares encontrados.}
  \leitura{linacre2022}{A descrição de uma biblioteca de código aberto de conciliação
    probabilística usada em escala por órgãos públicos.}
  \leitura{loshin2009}{Para os componentes de uma solução completa de master data, útil
    para saber o que se está deixando de fora ao começar pequeno.}
\end{leituras}
